% BHU PYQ -- Professional Exam Template
\documentclass[11pt,a4paper]{article}

\usepackage[a4paper,top=2.5cm,bottom=2.5cm,left=2.8cm,right=2.8cm,headheight=14pt]{geometry}
\usepackage{amsmath,amssymb}
\usepackage[T1]{fontenc}
\usepackage[utf8]{inputenc}
\usepackage{lmodern}
\usepackage{microtype}
\usepackage{fancyhdr}
\usepackage{enumitem}
\usepackage{listings}
\usepackage{hyperref}
\usepackage{lastpage}
\usepackage{parskip}

\hypersetup{colorlinks=true,urlcolor=black,linkcolor=black,
  pdftitle={BVCA-403: Cloud Computing -- BHU PYQ}}

\pagestyle{fancy}
\fancyhf{}
\renewcommand{\headrulewidth}{0.4pt}
\renewcommand{\footrulewidth}{0.4pt}
\fancyhead[L]{\small\textit{Banaras Hindu University}}
\fancyhead[R]{\small\textit{BVCA-403: Cloud Computing}}
\fancyfoot[C]{\small Page \thepage\ of \pageref{LastPage}}

\lstset{
  basicstyle=\small\ttfamily,
  frame=single,
  breaklines=true,
  columns=flexible,
  keepspaces=true
}

\newlist{parts}{enumerate}{2}
\setlist[parts,1]{label=(\alph*),leftmargin=*,itemsep=4pt,topsep=3pt}
\setlist[parts,2]{label=(\roman*),leftmargin=*,itemsep=2pt,topsep=2pt}
\newcommand{\pts}[1]{\hfill\textit{[#1]}}

\begin{document}
\begin{center}
  {\large\bfseries BANARAS HINDU UNIVERSITY}\\[2pt]
  {\normalsize B.Voc. Semester IV Examination, 2023-24}\\[6pt]
  \rule{0.8\linewidth}{0.5pt}\\[6pt]
  {\large\bfseries Paper: BVCA-403 --- Cloud Computing}\\[4pt]
  {\normalsize Time: Three Hours \hfill Maximum Marks: 70}
\end{center}

\rule{\linewidth}{0.4pt}

\smallskip
\noindent\textit{Instructions: This question paper carries three sections : Section |}

\bigskip

| 76993
: Paper : BVCA-403
| (Critical Design Thinking in Social Science)
| A, Band C. All questions are compulsory.
| SECTION -A
| Note: Answer any two of the following questions in
| 500 words approximately. Each question carries
| 12 marks.

\medskip\noindent\textbf{Question 1.} Explain the concept of critical thinking. Describe the
benefits of critical thinking.
)
: 2. Who is a critical thinker ? Discuss the characteristics of
a critical thinker.
| 3. Why design thinking skills matter ? How can we learn
design thinking ? Discuss.
Pie

\medskip\noindent\textbf{Question 76993.}

\medskip\noindent\textbf{Question 4.} What is the significance of studying logic ? What is not
an argument ? Explain.

\bigskip\noindent{\large\bfseries SECTION --- B |}
\rule{\linewidth}{0.4pt}\smallskip

Note: Answer any six of the following questions in about
200 words approximately. Each question carries
6 marks.

\medskip\noindent\textbf{Question 5.} (a) Explain the concept of sociocentrism.
\begin{parts}
  \item What is the importance of critical thinking at
  workplace ? Discuss.
  \item Describe the phases of design thinking.
  \item Explain fairness as critical thinking standard.
  (ec) What are the design thinking skills ? Describe
  team building as design thinking skill.
  \item Explain the functions of communication.
  \item Describe analytical skill as essential skill for
  design thinking.
  \item Differentiate between arguments and
  explanations.
  \item How can we identify premises and explanations ?
  Explain.
  SECTION --- C
  Note: Answer each of the following questions. Each
  question carries 1 mark.
\end{parts}

\medskip\noindent\textbf{Question 6.} (a) What do you mean by thinking ?
\begin{parts}
  \item Explain how one can improve communication
  skills.
  | 76993
  \item How logical correctness can be practiced ?
  \item What are the levels of communication ?
  | (e) List some of the most common barriers to critical
  | thinking.
  \item Explain the concept of logical correctness.
  | (g) What are the three mistakes which we make in
  thinking ?
(h) Elaborate the meaning of argument.
  \item Why analytical skills of important for a
  professional ?
  : (j) What are the 5 C's of effective communication ?
  !
  ) =) 200
  f eae nN i) I A ees Si;
  ‘
  ‘
  ' ‘
\end{parts}

\vfill
\rule{\linewidth}{0.4pt}
\begin{center}\small
\href{https://bhu-exam.vercel.app/}{Click here to visit our website}\\[4pt]\href{https://me-aryan.vercel.app/}{Click here to visit our contributor}\\[4pt]\href{https://quashan-me.vercel.app/}{Click here to visit our contributor}
\end{center}
\end{document}
